%
% menu_presence_locks_inverted_buggy.tex
%
% FIDELITY CONTROL for docs/menu_presence_locks.tex: the REJECTED design in
% which the registry fires the presence hook from inside its own lock. The
% only change from menu_presence_locks.tex is that the readers (BeginRead,
% EndRead) are replaced by BeginRegistryWrite and FireHook. A deadlock is
% reachable. The rest of this header is that of the fixed model.
%
% menu_presence_locks.tex
%
% Companion to docs/menu_presence.tex: the LOCK ORDER the fixed design
% requires, checked for deadlock-freedom.
%
% menu_presence.tex makes every presence-changing registry transition one
% atomic step. That atomicity is only real if the code can actually hold the
% locks it needs without a cycle. The locks in play:
%
%   StoreLock   the StoreLock every ClientLifecycle method takes (arrive,
%               identify, attach, renew, register_callback, detach, depart,
%               depart_lapsed), a re-entrant write lock; also held by the
%               Hub's scene writes and by ClientLifecycle.bind_sink. The
%               departure sinks run inside it.
%   RegLock     HubClientRegistry._lock, a plain (non-re-entrant) lock, which
%               PresenceWrites shares by reference. Taken INSIDE StoreLock by
%               every writer (PresenceWrites.record, renew_if_registered,
%               discard, reap_lapsed, attach,
%               register_callback, detach), and ALONE by the readers
%               (named_sessions, live_sessions, sessions, session_of, identity_of).
%   (marker)    The DirtyMarker's own lock, taken by mark_menus. It is a leaf:
%               nothing is acquired while it is held, so it contributes no
%               edge to the acquisition graph and is not modelled.
%
% Re-entrancy of StoreLock contributes no edge either: a thread already
% holding it never waits on it, so a nested acquire is a self-loop. It is
% modelled by omission.
%
% THE FIXED DESIGN: every presence-changing writer enters through
% ClientLifecycle (StoreLock first, then RegLock); the registry returns
% "presence changed" and the caller, still inside StoreLock, marks and binds.
% The callback route is no exception: ClientLifecycle.register_callback and
% ClientLifecycle.detach take the store lock, and only then reach
% PresenceWrites, which takes RegLock. (Before the fix CallbackOperations and
% ws_listen were handed the listener-leg write (then ListenerLegs, now PresenceWrites) and
% wrote under RegLock alone; that is
% deadlock-free, a registry-only writer being a reader as far as the lock
% graph goes, and it is wrong for the reason menu_presence.tex gives: no sink
% is bound.) Readers take RegLock alone and never wait for the store. The
% acquisition graph has the single edge StoreLock -> RegLock: acyclic, so
% deadlock-free.
%
% THE REJECTED DESIGN (menu_presence_locks_inverted_buggy.tex): the registry
% itself fires the mark/bind hook from inside record() (or, equally,
% PresenceWrites.register_callback), i.e. while holding RegLock, and binding a
% sink takes StoreLock. Writers that enter through the
% registry then hold RegLock and want StoreLock, against ClientLifecycle's writers
% that hold StoreLock and want RegLock: a cycle. ProB finds the deadlock.
%
% Type-check:  fuzz docs/menu_presence_locks.tex
% Model-check: probcli docs/menu_presence_locks.tex -model_check \
%                -p DEFAULT_SETSIZE 2
% Expected: no deadlock, no invariant violation (the fixed order).
%

\documentclass[a4paper,10pt,fleqn]{article}
\usepackage[margin=1in]{geometry}
\usepackage{fuzz}

\begin{document}

\title{lux Menu Presence: the Lock Order, inverted control}
\maketitle

\begin{zed}
  [T]
\end{zed}

A thread: an MCP tool thread, the listen leg's event-loop thread, the lease
sweep, the replicator tick. Two suffice for a cycle.

\begin{zed}
  PC ::= pidle | pstore | pboth | preg | pcstore | pcboth
\end{zed}

A thread's program counter: holding nothing; holding the store lock and
wanting the registry's; holding both; holding the registry lock alone (a
read); and the same two writing states for a thread on the callback route
(\texttt{pcstore}, \texttt{pcboth}), kept apart from the others so that the
model says in so many words that \texttt{register\_callback} and
\texttt{detach} follow the same order.

\begin{schema}{Locks}
  pc : T \fun PC \\
  storeHeld : \power T \\
  regHeld : \power T
\where
  \# storeHeld \leq 1 \\
  \# regHeld \leq 1 \\
  storeHeld = \{ t : T | pc~t = pstore \lor pc~t = pboth \lor
    pc~t = pcstore \lor pc~t = pcboth \} \\
  regHeld = \{ t : T | pc~t = pboth \lor pc~t = preg \lor
    pc~t = pcboth \}
\end{schema}

The lock tables are derived from the program counters, so they can never
disagree with them. The cardinality bounds are mutual exclusion, assumed as a
guard, as a lock is modelled: a thread cannot take a held lock. They are not
checked: what the model puts to the test is the order, by the deadlock check.

\begin{schema}{Init}
  Locks'
\where
  pc' = T \cross \{ pidle \} \\
  storeHeld' = \emptyset \\
  regHeld' = \emptyset
\end{schema}

\begin{schema}{BeginWrite}
  \Delta Locks \\
  t? : T
\where
  pc~t? = pidle \\
  storeHeld = \emptyset \\
  pc' = pc \oplus \{ t? \mapsto pstore \} \\
  storeHeld' = \{ t? \} \\
  regHeld' = regHeld
\end{schema}

A presence-changing writer takes the store lock first
(\texttt{ClientLifecycle.arrive}, \texttt{identify}, \texttt{attach},
\texttt{renew}, \texttt{depart}, \texttt{depart\_lapsed}, and the read loop,
which reaches the registry only through them).

\begin{schema}{TakeRegistry}
  \Delta Locks \\
  t? : T
\where
  pc~t? = pstore \\
  regHeld = \emptyset \\
  pc' = pc \oplus \{ t? \mapsto pboth \} \\
  storeHeld' = storeHeld \\
  regHeld' = \{ t? \}
\end{schema}

\begin{schema}{Finish}
  \Delta Locks \\
  t? : T
\where
  pc~t? = pboth \\
  pc' = pc \oplus \{ t? \mapsto pidle \} \\
  storeHeld' = \emptyset \\
  regHeld' = \emptyset
\end{schema}

The upsert or removal happens under both locks; the registry returns, and the
caller marks and binds under the store lock it still holds, then releases
both. (The marker's leaf lock, taken in between, is not modelled.)

\begin{schema}{BeginCallbackWrite}
  \Delta Locks \\
  t? : T
\where
  pc~t? = pidle \\
  storeHeld = \emptyset \\
  pc' = pc \oplus \{ t? \mapsto pcstore \} \\
  storeHeld' = \{ t? \} \\
  regHeld' = regHeld
\end{schema}

\begin{schema}{TakeRegistryForCallback}
  \Delta Locks \\
  t? : T
\where
  pc~t? = pcstore \\
  regHeld = \emptyset \\
  pc' = pc \oplus \{ t? \mapsto pcboth \} \\
  storeHeld' = storeHeld \\
  regHeld' = \{ t? \}
\end{schema}

\begin{schema}{FinishCallbackWrite}
  \Delta Locks \\
  t? : T
\where
  pc~t? = pcboth \\
  pc' = pc \oplus \{ t? \mapsto pidle \} \\
  storeHeld' = \emptyset \\
  regHeld' = \emptyset
\end{schema}

The callback route: \texttt{ClientLifecycle.register\_callback} and
\texttt{ClientLifecycle.detach} take the store lock, then
\texttt{PresenceWrites} takes the registry lock for its compare-and-write, then
the lifecycle marks and binds the one shared sink and releases both
(\texttt{RegisterCallback} and \texttt{Detach} in \texttt{menu\_presence.tex}).
The three steps have the shape of the three above, and the acquisition graph
gains no edge: StoreLock $\to$ RegLock, as before.

\begin{schema}{BeginRegistryWrite}
  \Delta Locks \\
  t? : T
\where
  pc~t? = pidle \\
  regHeld = \emptyset \\
  pc' = pc \oplus \{ t? \mapsto preg \} \\
  storeHeld' = storeHeld \\
  regHeld' = \{ t? \}
\end{schema}

A writer that enters through the registry directly, as
\texttt{ws\_listen}'s read loop and \texttt{CallbackOperations} did when they
were handed the registry and the listener-leg write (then
\texttt{ListenerLegs}, now \texttt{PresenceWrites})
(\texttt{clients.record}, \texttt{legs.register\_callback}).

\begin{schema}{FireHook}
  \Delta Locks \\
  t? : T
\where
  pc~t? = preg \\
  storeHeld = \emptyset \\
  pc' = pc \oplus \{ t? \mapsto pboth \} \\
  storeHeld' = \{ t? \} \\
  regHeld' = regHeld
\end{schema}

The defect: the registry (or \texttt{PresenceWrites}) fires the mark/bind hook
from inside \texttt{record} or \texttt{register\_callback}, while holding its
own lock, and binding a sink takes the store lock. The edge RegLock $\to$
StoreLock closes a cycle with BeginWrite's StoreLock $\to$ RegLock.

Expected: a DEADLOCK is found.

\begin{verbatim}
  probcli docs/menu_presence_locks_inverted_buggy.tex -model_check \
    -p DEFAULT_SETSIZE 2
\end{verbatim}

Trace: BeginWrite(t1); BeginRegistryWrite(t2) leaves t1 waiting for the
registry (held by t2) and t2 waiting for the store (held by t1).

\end{document}
